\clearpage\fancyhead[L]{\small GAUSS TO VON LINDENAU.}\hypertarget{native-331-lower}{}\label{native-331-lower}
\begin{center}\textbf{GAUSS TO VON LINDENAU.}\par\vspace{4pt}\hrule\vspace{4pt}\fontsize{10}{12}\selectfont Monatliche Correspondenz. Vol. XXIV. pp. 180\,..\,182. August 1811.\par\vspace{4pt}\hrule\vspace{4pt}\end{center}

\begin{flushright}Göttingen, 3 August 1811.\end{flushright}
\Needspace{16\baselineskip}\par\indent In the last few days I have provisionally derived the following parabolic elements from Mr. \textsc{von Zach}’s observations of the comet, which you had sent to Professor \textsc{Harding}.
\SourceLeaf{332}
\begin{center}\fontsize{10}{12}\selectfont Perihelion passage 1811 Sept. 10\\$0^{\text{u}}\;51^{\text{m}}\;21^{\text{s}}$ Göttingen meridian.\end{center}

\begin{center}\fontsize{10}{12}\selectfont\setlength{\tabcolsep}{4pt}\renewcommand{\arraystretch}{1.15}\begin{tabular}{llr}
Logarithm of the distance at perihelion & .\;.\;.\;.\;.\;.\;. & $9.99153$\\
Longitude of perihelion & .\;.\;.\;.\;.\;.\;.\;.\;.\;.\;.\;.\;.\;.\;. & $73^\circ\;14'\;35''$\\
Longitude of the ascending node & .\;.\;.\;.\;.\;.\;.\;.\;.\;.\;.\;. & $141\;4\;59$\\
Inclination of the orbit & .\;.\;.\;.\;.\;.\;.\;.\;.\;.\;.\;.\;.\;.\;.\;.\;. & $73\;48\;2$\\
Retrograde motion &  & \\
\end{tabular}\end{center}

\Needspace{15\baselineskip}\par\indent After the middle of this month one can already start looking for this remarkable celestial body in the northern sky. For that purpose I have calculated the following positions from the elements above.

\begin{center}\fontsize{10}{12}\selectfont\setlength{\tabcolsep}{4pt}\renewcommand{\arraystretch}{1.15}\begin{tabular}{lrr}
1811 Aug. 11 $12^{\text{u}}$ & AR. $142^\circ\;43'$ & Decl. $29^\circ\;54'$ N.\\
26 & $152\;19$ & $36\;15$\\
Sept. 10 & $167\;34$ & $43\;14$\\
\end{tabular}\end{center}

\begin{center}\fontsize{10}{12}\selectfont\setlength{\tabcolsep}{4pt}\renewcommand{\arraystretch}{1.15}\begin{tabular}{lrr}
1811 Sept. 25 $12^{\text{u}}$ & AR. $194^\circ\;22'$ & Decl. $48^\circ\;35'$ N.\\
Oct. 10 & $232\;28$ & $44\;52$\\
25 & $262\;49$ & $30\;42$\\
\end{tabular}\end{center}

\Needspace{20\baselineskip}\par\indent These results can only be regarded as provisional, however, since the slow heliocentric motion so far (only $15^\circ$ from 11 April to 2 June) does not yet allow a very precise determination of the elements. The brightness, taking that at distance 1 from Earth and Sun as the unit, was as follows

\begin{center}\fontsize{10}{12}\selectfont\setlength{\tabcolsep}{4pt}\renewcommand{\arraystretch}{1.15}\begin{tabular}{llr}
11 April & .\;.\;.\;.\;. & $0.035$\\
7 May & .\;.\;.\;.\;. & $0.042$\\
2 June & .\;.\;.\;.\;. & $0.052$\\
\end{tabular}\end{center}

\noindent and will be as follows

\begin{center}\fontsize{10}{12}\selectfont\setlength{\tabcolsep}{4pt}\renewcommand{\arraystretch}{1.15}\begin{tabular}{llr}
Aug. 11 & .\;.\;.\;.\;.\;.\;.\;.\;. & $0.194$\\
26 & .\;.\;.\;.\;.\;.\;.\;.\;. & $0.294$\\
Sept. 10 & .\;.\;.\;.\;.\;.\;.\;.\;. & $0.434$\\
Sept. 25 & .\;.\;.\;.\;.\;.\;.\;.\;. & $0.576$\\
Oct. 10 & .\;.\;.\;.\;.\;.\;.\;.\;. & $0.595$\\
25 & .\;.\;.\;.\;.\;.\;.\;.\;. & $0.415$\\
\end{tabular}\end{center}

\noindent Only around February will it become invisible in Aquarius. It may become visible there again in April or May, if its light is not too faint. For 2 July 1812 I made the following rough calculation.
\begin{center}R.A. $336^\circ$\quad Decl. $22^\circ$ S\quad Brightness $0.005$.\end{center}
\par\indent [Handwritten remark. From improved elements 1812 June $3.22^{\mathrm u}$ $337^\circ\,49'$ --- $11^\circ\,41'$ --- $0.5435$ $0.0057$.]

\par\indent I have now continued the calculations of the perturbations of \emph{Pallas} to the end of 1812. In a few days I hope to send you the ephemeris, which Mr. \textsc{Nicolai}, a very experienced young man, is now calculating.
